\documentclass{article}
% ICLR template (2026 style files as base; swap to iclr2027_conference.* on release)
\usepackage{iclr2026_conference,times}
\usepackage{natbib}
\usepackage{amsmath,amssymb}
\usepackage{booktabs}
\usepackage{graphicx}
\usepackage{hyperref}
\usepackage{url}
\usepackage{algorithm}
\usepackage{algorithmic}

% \iclrfinalcopy  % uncomment for camera-ready (de-anonymizes)

\title{Wiring Beats Blending: What Transfers Between\\Transformer Sizes --- and What Doesn't}

\author{Anonymous authors\\
\texttt{(double-blind submission)}}

\begin{document}
\maketitle

\begin{abstract}
Model families train every size from scratch. Can a pretrained large model be
converted into a smaller sibling? We characterize the 1.4B$\to$410M conversion
in the Pythia family end-to-end: (i)~representations align strongly across
sizes (ridge $R^2=0.84$) while parameters align weakly; (ii)~dense weight
projection is functionally destructive --- provably not an assembly artifact ---
because basis mixing breaks rotary, per-head, GELU, and LayerNorm structure;
(iii)~after the best-fit linear operator, weight residuals are statistically
indistinguishable from noise under shuffle controls; (iv)~conversion value
therefore lives in \emph{initialization}: in matched-budget continued
pre-training we decompose conversion into two independent levers ---
least-squares \textbf{compensation} (function: best zero-shot) and
variance-preserving \textbf{rescale} (training dynamics: best endpoints) ---
whose combination dominates the strongest subcloning variant on both axes and
in every seed ($83.9\pm1.8$ vs.\ $89.7\pm3.7$ final perplexity, 3/3 paired
wins; 18$\times$ better than from-scratch at 30M tokens). The ordering holds
out-of-domain and at $2\times$ the training context, and the $\sim$13$\times$
margin over from-scratch generalizes with three seeds to a held-out,
depth-dominated pair --- where selection variants tie, mapping the boundary of
compensation to the width-reduced regime. Code, checkpoints, and the frozen
evaluation corpus are released.
\end{abstract}

\section{Introduction}
% WRITE LAST (per plan). Funnel: cost motivation -> italicized research question
% -> approach -> 4 contribution bullets (see project_docs/reports/paper1_outline.md).

\section{Related work}
\label{sec:related}

\paragraph{Downscaling pretrained models.}
The closest prior work is Weight Subcloning~\citep{samragh2023weight}, which
initializes a smaller transformer by ranking units of a larger pretrained one,
copying selected rows and columns, and rescaling, establishing that
selection-based initialization accelerates training of the target size. Sheared
LLaMA~\citep{xia2024sheared} learns structured pruning masks jointly with
continued pre-training at a fraction of from-scratch compute; it optimizes the
full pipeline, whereas we hold the training recipe fixed across arms and isolate
what the initialization alone contributes. Our closed-form compensation
belongs to the pruning-with-reconstruction lineage (adjusting surviving
weights to absorb the function of removed ones, from Optimal Brain
Surgeon~\citep{hassibi1993optimal} to structured LLM pruning like
LLM-Pruner~\citep{ma2023llmpruner}), applied once, in closed form, at conversion
time. Knowledge distillation~\citep{hinton2015distilling} transfers function
through a teacher's outputs and requires a full training run; we study the
complementary question of what transfers through weights at negligible
cost, and the two are combinable.

\paragraph{Growing pretrained models.}
The mirror direction (initializing a larger model from a smaller one)
has a longer history: Net2Net~\citep{chen2016net2net} introduced
function-preserving widening and deepening, bert2BERT~\citep{chen2022bert2bert}
adapted it to transformers, staged training~\citep{shen2022staged} formalized
growth operators that preserve loss and dynamics, and LiGO~\citep{wang2023ligo}
learns a structured linear map from small weights to the large initialization.
That linear operators carry useful signal upward is consistent with our finding
that, downward, even a fitted dense projection initializes far better than random;
the asymmetry we document is that dense basis mixing destroys the rotary,
per-head, and normalization structure that selection preserves
(\S\ref{sec:projection}).

\paragraph{Representation similarity and parameter symmetries.}
Our analyses use standard similarity tools:
CKA~\citep{kornblith2019similarity} and SVCCA~\citep{raghu2017svcca} measure
representation alignment, and model
stitching~\citep{lenc2015understanding,bansal2021revisiting} tests functional
interchangeability through a trained adapter. Work on permutation symmetries and
model merging~\citep{ainsworth2023git} shows that networks are naturally compared
modulo the transformations under which the architecture is invariant. Our
selection-versus-blending result is an instance of the same principle: structured
selection composes a permutation (an element of the architecture's symmetry
group, restricted to respect head boundaries and rotary frequency
pairs~\citep{su2024roformer}) with coordinate deletion, so every surviving unit
keeps the exact nonlinear and positional semantics the architecture assigns it;
dense projection mixes coordinates, exits the group, and is punished for it.

\paragraph{Positioning.}
Each ingredient above exists in isolation: selection-based initialization,
least-squares reconstruction, linear growth operators, representation similarity.
What is new is the setting (conversion between sizes of a single family
trained on identical data, removing data and tokenizer confounds) together with
rotary-frequency-matched selection and a controlled end-to-end characterization,
from representation alignment through projection diagnosis and residual null tests
to a matched-budget race. Because our subcloning baseline re-implements the recipe
of \citet{samragh2023weight}, we disclose three fidelity differences and label it
\emph{subcloning-style}: (a)~we score attention heads and MLP units by weight
norms and residual lanes by activation variance, where they score by activation
magnitudes; (b)~when depth must shrink we remove blocks at even stride rather than
from the middle; and (c)~we evaluate their $\sqrt{d/d'}$ weight rescale directly
rather than adopting it wholesale (\S\ref{sec:race}): applied to every matrix it
collapses zero-shot quality because LayerNorm already renormalizes most read
paths, while its variance-preserving motivation holds exactly on the projections
no norm protects, the same two families our compensation re-fits.

\section{Experimental setup}
\label{sec:setup}

\paragraph{Model family and conversion pairs.}
We study the Pythia suite~\citep{biderman2023pythia}: a family of GPT-NeoX
models trained on identical data (the Pile) with a shared tokenizer, which
isolates \emph{size} as the only variable between family members. Our primary
pair converts \textbf{1.4B}$\,\to\,$\textbf{410M} (24 layers and 16 heads in
both; only widths shrink: residual $2048\!\to\!1024$, head dim
$128\!\to\!64$, MLP $8192\!\to\!4096$), so layer correspondence is trivially
one-to-one and the analysis isolates width. A held-out pair,
\textbf{410M}$\,\to\,$\textbf{160M}, exercises the axes the primary pair does
not: depth halves ($24\!\to\!12$ blocks), the head \emph{count} drops
($16\!\to\!12$; head dim unchanged), and width cuts are mild
($\sim$25\%). A third pair, \textbf{6.9B}$\,\to\,$\textbf{1.4B}, tests scale
(all three axes reduced; \S\ref{sec:race}).
% TODO(owner): confirm final pair list once scale runs land.

\paragraph{Frozen evaluation corpus.}
All representation analyses probe every model with the \emph{same} inputs:
10{,}000 sequences of 128 tokens drawn from the Pile
(\texttt{pile-uncopyrighted}), tokenized once and frozen for the entire
project. The first 500 sequences (all token positions) form the subset for
similarity analysis (\S\ref{sec:representations}); the remainder contribute
mean-pooled and last-token activations. Continued pre-training streams the
same corpus but skips the region used to build the frozen set (a
contamination guard), and each random seed trains on a disjoint stream
offset.

\paragraph{Behavioral evaluation.}
Language-modeling quality is strided sliding-window perplexity on
WikiText-103 validation (every token scored exactly once). Hardening
evaluations add C4 perplexity (out-of-domain), perplexity at $2\times$ the
training context (2048), and zero-shot accuracy on LAMBADA, ARC-Easy,
HellaSwag, and PIQA via \texttt{lm-eval}. Reference models anchor every
table; anchor scores match published Pythia numbers.

\paragraph{Determinism and verification.}
All forward passes are greedy and seeded; repeated runs are bit-identical on
the same device. Weight extraction is validated by exact reconstruction: the
fused GPT-NeoX QKV matrices are split head-aware and re-fused to bit-identical
equality, and a model rebuilt from its own extracted weights reproduces the
original logits with zero difference (\S\ref{sec:projection}) --- any quality
change downstream is therefore attributable to conversion mathematics, not
assembly.

\paragraph{Training protocol.}
Every conversion arm within an experiment shares \emph{identical} data order,
optimizer (AdamW, $\beta=(0.9,0.95)$, weight decay $0.1$), cosine schedule
with 100-step warmup, gradient clipping at $1.0$, bf16 autocast, and token
budget (30M primary; 100M persistence checks); the learning rate equals the
target size's original pretraining rate. Initialization construction costs
are negligible against any training budget (selection: seconds; moments pass
for compensation: $\sim$75\,s; closed-form solves: seconds) and are included
in the compute accounting. All experiments ran on a single NVIDIA GB10
(128\,GB unified memory); any 24\,GB+ CUDA GPU reproduces them
(Appendix~TODO).

\section{Representations align, parameters do not}
% From M2.md: CKA heatmap (saturated middle band), ridge 0.84 vs Procrustes
% 0.72, embeddings weight-space R2 0.25-0.39. Figures: cka_heatmap, maps_r2.

\section{Why dense projection fails}
% From M3.md: target-free projection residuals >1; ppl anchors incl. LN
% correction (181k -> 12.3k); identity-reconstruction bit-exact; diagnosis:
% rotary / per-head / GELU / per-feature-LN break under basis mixing.

\section{Post-operator residuals are noise}
\label{sec:residuals}

Section~\ref{sec:projection} shows that even the strongest shared linear operator
explains at best $\sim$30\% of weight variance (relative error $0.66$--$0.80$),
leaving the majority as a residual. Is that residual a structured, learnable
correction, or noise? We define the residual per weight as
\begin{equation}
\Delta(l,\text{type}) \;=\; W_{410\text{M}} \;-\; A\, W_{1.4\text{B}}\, B^{\top},
\label{eq:delta}
\end{equation}
using the shared fitted operator $(A,B)$ per type, the strongest-alignment
condition of \S\ref{sec:projection}, so whatever remains is what no single linear
size-map can explain and the verdict is conservative: if any signal survives the
best linear map, these tests should find it. Four tests probe $\Delta$, each
against a matched control.

\paragraph{Test 1: spectral (effective rank).}
App.~Table~\ref{tab:erank} compares the effective rank of $\Delta$ against a
per-type shuffled control and a shape/scale-matched Gaussian control. $\Delta$
sits $2.8$--$6.4\%$ below both controls in every type: a faint spectral
concentration, so the residual is not perfectly isotropic noise
(App.~Fig.~\ref{fig:spectra}).

\paragraph{Test 2: cross-layer consistency.}
The mean pairwise cosine of vectorized $\Delta$ across layers is
$\approx\pm0.0002$ in every type, identical to the shuffle control: no direction
is shared across layers, so no single correction is uniformly missing.

\paragraph{Test 3: learnability (decisive).}
A patch predictor $R(\hat{W}\text{ patch} + \text{type/layer/position})
\to \Delta\text{ patch}$ is trained on 18 layers, scored on 6 stratified held-out
layers as error reduction over predicting zero; the control retrains on $\Delta$
shuffled across layers, with a 1000-sample bootstrap on the real$-$control gap.
Neither a linear nor an MLP-512 predictor beats zero (reductions
$-0.0004$/$-0.0003$), and real and control are indistinguishable (gap
CI$_{95}\approx[-0.0000,-0.0000]$, $p\ge0.99$; Table~\ref{tab:predictor}). No
patch-learnable signal.

\paragraph{Test 4: transfer.}
On the held-out 410M$\,\to\,$160M pair the shared operators fit at relative error
$0.56$--$0.74$ (smaller width gap). A predictor trained on all primary-pair
layers, evaluated on the held-out pair, reduces error by $-0.0004$/$-0.0003$,
again indistinguishable from its shuffle control ($p=1.00$;
Table~\ref{tab:predictor}). Nothing transfers, consistent with Test~3.

\paragraph{Reconciling the faint spectral signal.}
Tests 1 and 3 are not in tension: the $2.8$--$6.4\%$ spectral deficit is a
generic statistical trace, not a predictor-exploitable correspondence, and
a behavioral cross-check confirms the destruction lives in the dense projection,
not the residual (Appendix~\ref{app:reconcile}).

\paragraph{Handoff to \S\ref{sec:race}.}
Within a family the sizes share representation geometry but not parameter
solutions, and what the best linear size-map misses behaves as noise (no
cross-layer direction, no patch-learnable mapping, no cross-pair transfer; only a
faint generic spectral fingerprint), so conversion value must come from
initialization plus brief fine-tuning, the matched-budget recovery race
of \S\ref{sec:race}.

\ificlrcompact\else\TabPredictor\fi  % deferred to appendix in the ICLR build

\section{Matched-budget conversion: the two levers}
% From M5.md: 7-arm ablation ladder, two-lever decomposition (compensation =
% function, rescale = training dynamics), seeded verdict 83.9+-1.8 vs
% 89.7+-3.7 (3/3 paired), hardening (6 metrics / 100M / unseen pair 3 seeds),
% [SCALE TABLE 6.9B->1.4B pending tonight's runs]. Algorithm box: selection
% scoring + ls_compensate. Compute accounting.

\section{Limitations and future work}
% Single family (port to Llama/Qwen/DeepSeek: per-arch safe-wiring analysis);
% budgets << full recovery; seeds; read-in compensation unexplored;
% reduction-aware conversion (width-vs-depth allocation, block distill).


\subsubsection*{Reproducibility statement}
% TODO: repo + HF links (public at camera-ready), uv.lock pinning, determinism
% notes (bit-identical per-GPU), frozen-corpus recipe, all seeds/configs in repo.

\bibliography{references}
\bibliographystyle{iclr2026_conference}

\end{document}
